% =====================================================================
%  UI EVIDENCE FRAMES  --  % =====================================================================
%  UI EVIDENCE FRAMES  --  % =====================================================================
%  UI EVIDENCE FRAMES  --  % =====================================================================
%  UI EVIDENCE FRAMES  --  \input{frames_ui} from the main .tex
%
%  Images: search.png, similar.png -- captured from the running frontend
%  on 14 Sep 2026 (Vite dev server, profile "Sci-Fi Explorer", 1600x1000).
%  Every number quoted below is read off those two captures. If you
%  re-shoot the UI, re-read the numbers: they WILL change, and a panel
%  that checks a sub-score against the slide is the worst way to find out.
%
%  HEIGHT BUDGET: a 16:9 frame gives ~6.8cm of body. The screenshot
%  heights below are what keeps the bullets on the slide -- if you
%  enlarge an image, delete a bullet to pay for it.
% =====================================================================

% ---------------------------------------------------------------------
\begin{frame}[shrink]{2. System in Action --- Search and the Score Breakdown}
\framesubtitle{Every ranked card exposes the sub-scores that produced it}
\vspace{-1mm}
\begin{center}
\shot{search}{width=0.84\linewidth,height=4.0cm,keepaspectratio}%
{Search results with one card's score breakdown expanded.}
\end{center}
\vspace{-1mm}
\begin{columns}[T]
\column{0.485\textwidth}
\blk{What is on screen}
\begin{itemize}\scriptsize\setlength{\itemsep}{2pt}
  \item \texttt{"space adventure with aliens"}, profile \texttt{Sci-Fi Explorer} --- \textbf{10
        results in 3973\,ms}, spanning movies and shows.
  \item \textbf{\#1 Aliens in the Attic, 0.912.} The bars are the score, not decoration.
\end{itemize}

\column{0.485\textwidth}
\blk{The sub-scores sum to the printed total}
\begin{itemize}\scriptsize\setlength{\itemsep}{2pt}
  \item $0.681 + 0.100 + 0.090 + 0.040 = \mathbf{0.911}$ against a printed \textbf{0.912} ---
        rounding, and checkable on screen.
  \item \textcolor{amritaMaroon}{\textbf{Graph is struck through}}: no LightGCN embedding, so the
        other four still rank it. Graceful degradation, live.
\end{itemize}
\end{columns}
\end{frame}

% ---------------------------------------------------------------------
\begin{frame}[shrink]{2. System in Action --- Find Similar and the Feedback Loop}
\framesubtitle{Item-to-item retrieval through \texttt{POST /recommend/item} --- same heads, different anchor}
\vspace{0.5mm}
\begin{columns}[T]
\column{0.45\textwidth}
\vspace{-1mm}
\begin{center}
\shot{similar}{width=\linewidth,height=5.2cm,keepaspectratio}%
{The Find-similar drawer, anchored on a result card.}
\end{center}

\column{0.525\textwidth}
\blk{Same five heads, a different anchor}
\begin{itemize}\scriptsize\setlength{\itemsep}{3pt}
  \item The query vector \emph{is} the anchor item's own 384-d embedding --- nothing is typed.
  \item \textbf{\#1 Minions \& Monsters, 0.847}
        $= 0.535 + 0.128 + 0.101 + 0.063 + 0.021$.
  \item \textcolor{amritaMaroon}{\textbf{Graph contributes $+0.128$ here}}, where it was struck out
        on the previous slide --- the \textbf{0.20\%} artifact coverage, visible in the product.
\end{itemize}
\vspace{1mm}
\blk{Closing the loop}
\begin{itemize}\scriptsize\setlength{\itemsep}{3pt}
  \item like / skip $\rightarrow$ \texttt{/interactions} $\rightarrow$ MySQL $\rightarrow$ EMA
        ($\alpha = 0.25$, \textbf{39\,$\mu$s}).
  \item The next query ranks from the updated vector --- never a retrain.
\end{itemize}
\end{columns}
\end{frame}
 from the main .tex
%
%  Images: search.png, similar.png -- captured from the running frontend
%  on 14 Sep 2026 (Vite dev server, profile "Sci-Fi Explorer", 1600x1000).
%  Every number quoted below is read off those two captures. If you
%  re-shoot the UI, re-read the numbers: they WILL change, and a panel
%  that checks a sub-score against the slide is the worst way to find out.
%
%  HEIGHT BUDGET: a 16:9 frame gives ~6.8cm of body. The screenshot
%  heights below are what keeps the bullets on the slide -- if you
%  enlarge an image, delete a bullet to pay for it.
% =====================================================================

% ---------------------------------------------------------------------
\begin{frame}[shrink]{2. System in Action --- Search and the Score Breakdown}
\framesubtitle{Every ranked card exposes the sub-scores that produced it}
\vspace{-1mm}
\begin{center}
\shot{search}{width=0.84\linewidth,height=4.0cm,keepaspectratio}%
{Search results with one card's score breakdown expanded.}
\end{center}
\vspace{-1mm}
\begin{columns}[T]
\column{0.485\textwidth}
\blk{What is on screen}
\begin{itemize}\scriptsize\setlength{\itemsep}{2pt}
  \item \texttt{"space adventure with aliens"}, profile \texttt{Sci-Fi Explorer} --- \textbf{10
        results in 3973\,ms}, spanning movies and shows.
  \item \textbf{\#1 Aliens in the Attic, 0.912.} The bars are the score, not decoration.
\end{itemize}

\column{0.485\textwidth}
\blk{The sub-scores sum to the printed total}
\begin{itemize}\scriptsize\setlength{\itemsep}{2pt}
  \item $0.681 + 0.100 + 0.090 + 0.040 = \mathbf{0.911}$ against a printed \textbf{0.912} ---
        rounding, and checkable on screen.
  \item \textcolor{amritaMaroon}{\textbf{Graph is struck through}}: no LightGCN embedding, so the
        other four still rank it. Graceful degradation, live.
\end{itemize}
\end{columns}
\end{frame}

% ---------------------------------------------------------------------
\begin{frame}[shrink]{2. System in Action --- Find Similar and the Feedback Loop}
\framesubtitle{Item-to-item retrieval through \texttt{POST /recommend/item} --- same heads, different anchor}
\vspace{0.5mm}
\begin{columns}[T]
\column{0.45\textwidth}
\vspace{-1mm}
\begin{center}
\shot{similar}{width=\linewidth,height=5.2cm,keepaspectratio}%
{The Find-similar drawer, anchored on a result card.}
\end{center}

\column{0.525\textwidth}
\blk{Same five heads, a different anchor}
\begin{itemize}\scriptsize\setlength{\itemsep}{3pt}
  \item The query vector \emph{is} the anchor item's own 384-d embedding --- nothing is typed.
  \item \textbf{\#1 Minions \& Monsters, 0.847}
        $= 0.535 + 0.128 + 0.101 + 0.063 + 0.021$.
  \item \textcolor{amritaMaroon}{\textbf{Graph contributes $+0.128$ here}}, where it was struck out
        on the previous slide --- the \textbf{0.20\%} artifact coverage, visible in the product.
\end{itemize}
\vspace{1mm}
\blk{Closing the loop}
\begin{itemize}\scriptsize\setlength{\itemsep}{3pt}
  \item like / skip $\rightarrow$ \texttt{/interactions} $\rightarrow$ MySQL $\rightarrow$ EMA
        ($\alpha = 0.25$, \textbf{39\,$\mu$s}).
  \item The next query ranks from the updated vector --- never a retrain.
\end{itemize}
\end{columns}
\end{frame}
 from the main .tex
%
%  Images: search.png, similar.png -- captured from the running frontend
%  on 14 Sep 2026 (Vite dev server, profile "Sci-Fi Explorer", 1600x1000).
%  Every number quoted below is read off those two captures. If you
%  re-shoot the UI, re-read the numbers: they WILL change, and a panel
%  that checks a sub-score against the slide is the worst way to find out.
%
%  HEIGHT BUDGET: a 16:9 frame gives ~6.8cm of body. The screenshot
%  heights below are what keeps the bullets on the slide -- if you
%  enlarge an image, delete a bullet to pay for it.
% =====================================================================

% ---------------------------------------------------------------------
\begin{frame}[shrink]{2. System in Action --- Search and the Score Breakdown}
\framesubtitle{Every ranked card exposes the sub-scores that produced it}
\vspace{-1mm}
\begin{center}
\shot{search}{width=0.84\linewidth,height=4.0cm,keepaspectratio}%
{Search results with one card's score breakdown expanded.}
\end{center}
\vspace{-1mm}
\begin{columns}[T]
\column{0.485\textwidth}
\blk{What is on screen}
\begin{itemize}\scriptsize\setlength{\itemsep}{2pt}
  \item \texttt{"space adventure with aliens"}, profile \texttt{Sci-Fi Explorer} --- \textbf{10
        results in 3973\,ms}, spanning movies and shows.
  \item \textbf{\#1 Aliens in the Attic, 0.912.} The bars are the score, not decoration.
\end{itemize}

\column{0.485\textwidth}
\blk{The sub-scores sum to the printed total}
\begin{itemize}\scriptsize\setlength{\itemsep}{2pt}
  \item $0.681 + 0.100 + 0.090 + 0.040 = \mathbf{0.911}$ against a printed \textbf{0.912} ---
        rounding, and checkable on screen.
  \item \textcolor{amritaMaroon}{\textbf{Graph is struck through}}: no LightGCN embedding, so the
        other four still rank it. Graceful degradation, live.
\end{itemize}
\end{columns}
\end{frame}

% ---------------------------------------------------------------------
\begin{frame}[shrink]{2. System in Action --- Find Similar and the Feedback Loop}
\framesubtitle{Item-to-item retrieval through \texttt{POST /recommend/item} --- same heads, different anchor}
\vspace{0.5mm}
\begin{columns}[T]
\column{0.45\textwidth}
\vspace{-1mm}
\begin{center}
\shot{similar}{width=\linewidth,height=5.2cm,keepaspectratio}%
{The Find-similar drawer, anchored on a result card.}
\end{center}

\column{0.525\textwidth}
\blk{Same five heads, a different anchor}
\begin{itemize}\scriptsize\setlength{\itemsep}{3pt}
  \item The query vector \emph{is} the anchor item's own 384-d embedding --- nothing is typed.
  \item \textbf{\#1 Minions \& Monsters, 0.847}
        $= 0.535 + 0.128 + 0.101 + 0.063 + 0.021$.
  \item \textcolor{amritaMaroon}{\textbf{Graph contributes $+0.128$ here}}, where it was struck out
        on the previous slide --- the \textbf{0.20\%} artifact coverage, visible in the product.
\end{itemize}
\vspace{1mm}
\blk{Closing the loop}
\begin{itemize}\scriptsize\setlength{\itemsep}{3pt}
  \item like / skip $\rightarrow$ \texttt{/interactions} $\rightarrow$ MySQL $\rightarrow$ EMA
        ($\alpha = 0.25$, \textbf{39\,$\mu$s}).
  \item The next query ranks from the updated vector --- never a retrain.
\end{itemize}
\end{columns}
\end{frame}
 from the main .tex
%
%  Images: search.png, similar.png -- captured from the running frontend
%  on 14 Sep 2026 (Vite dev server, profile "Sci-Fi Explorer", 1600x1000).
%  Every number quoted below is read off those two captures. If you
%  re-shoot the UI, re-read the numbers: they WILL change, and a panel
%  that checks a sub-score against the slide is the worst way to find out.
%
%  HEIGHT BUDGET: a 16:9 frame gives ~6.8cm of body. The screenshot
%  heights below are what keeps the bullets on the slide -- if you
%  enlarge an image, delete a bullet to pay for it.
% =====================================================================

% ---------------------------------------------------------------------
\begin{frame}[shrink]{2. System in Action --- Search and the Score Breakdown}
\framesubtitle{Every ranked card exposes the sub-scores that produced it}
\vspace{-1mm}
\begin{center}
\shot{search}{width=0.84\linewidth,height=4.0cm,keepaspectratio}%
{Search results with one card's score breakdown expanded.}
\end{center}
\vspace{-1mm}
\begin{columns}[T]
\column{0.485\textwidth}
\blk{What is on screen}
\begin{itemize}\scriptsize\setlength{\itemsep}{2pt}
  \item \texttt{"space adventure with aliens"}, profile \texttt{Sci-Fi Explorer} --- \textbf{10
        results in 3973\,ms}, spanning movies and shows.
  \item \textbf{\#1 Aliens in the Attic, 0.912.} The bars are the score, not decoration.
\end{itemize}

\column{0.485\textwidth}
\blk{The sub-scores sum to the printed total}
\begin{itemize}\scriptsize\setlength{\itemsep}{2pt}
  \item $0.681 + 0.100 + 0.090 + 0.040 = \mathbf{0.911}$ against a printed \textbf{0.912} ---
        rounding, and checkable on screen.
  \item \textcolor{amritaMaroon}{\textbf{Graph is struck through}}: no LightGCN embedding, so the
        other four still rank it. Graceful degradation, live.
\end{itemize}
\end{columns}
\end{frame}

% ---------------------------------------------------------------------
\begin{frame}[shrink]{2. System in Action --- Find Similar and the Feedback Loop}
\framesubtitle{Item-to-item retrieval through \texttt{POST /recommend/item} --- same heads, different anchor}
\vspace{0.5mm}
\begin{columns}[T]
\column{0.45\textwidth}
\vspace{-1mm}
\begin{center}
\shot{similar}{width=\linewidth,height=5.2cm,keepaspectratio}%
{The Find-similar drawer, anchored on a result card.}
\end{center}

\column{0.525\textwidth}
\blk{Same five heads, a different anchor}
\begin{itemize}\scriptsize\setlength{\itemsep}{3pt}
  \item The query vector \emph{is} the anchor item's own 384-d embedding --- nothing is typed.
  \item \textbf{\#1 Minions \& Monsters, 0.847}
        $= 0.535 + 0.128 + 0.101 + 0.063 + 0.021$.
  \item \textcolor{amritaMaroon}{\textbf{Graph contributes $+0.128$ here}}, where it was struck out
        on the previous slide --- the \textbf{0.20\%} artifact coverage, visible in the product.
\end{itemize}
\vspace{1mm}
\blk{Closing the loop}
\begin{itemize}\scriptsize\setlength{\itemsep}{3pt}
  \item like / skip $\rightarrow$ \texttt{/interactions} $\rightarrow$ MySQL $\rightarrow$ EMA
        ($\alpha = 0.25$, \textbf{39\,$\mu$s}).
  \item The next query ranks from the updated vector --- never a retrain.
\end{itemize}
\end{columns}
\end{frame}
